% =============================================================================
%  LEVITATION LAB — TECHNICAL DOCUMENTATION
%  GT4Pebbles Project
% =============================================================================
%
%  Compilation:
%    pdflatex -shell-escape levitation_lab_technical.tex
%
%  CHANGELOG (checked against source, this revision):
%    - §12: corrected tray sweep angle 270° → 240° (Math.PI * 4/3)
%    - §12: clarified that TUNING.tray.brakeDur is unused; actual brake
%            duration is computed dynamically from drum speed at fire time
%    - §17: corrected snap sound upper frequency 4300 Hz → 4100 Hz
%    - §17: corrected injection puff lifetime 0.42 s → 1.42 s
%    - §17: added soundTrayBeep entry (previously undocumented)
%    - §14: added missing DT values for game levels 3 and 5
%    - §15: added ?challenge&time=N URL parameter
%    - §11.3: encounter overlay updated — tunable showPairs/arcPairs,
%             three-tier colour scheme, orbit arcs, subtler flash
%    - §12: inHighlightSince cleared on tray capture; onTray guard added
%             to all levitation functions and encounter cache
%    - §13 (Reset): Reset now disables zoom mode
%    - Appendix A: added encounters.* tuning section; added tray.*,
%                  ball.bumpStrength, ball.settleVel, ball.contactDrag
%    - §10.4: simulation speed gears extended to /64× -- 8× (10 gears)
%    - §11.7: new subsection — live KDE orbit band (btnVtProj mode 2)
%    - §14 Expert HUD: updated btnVtProj (three-state) and
%             btnAggEncounters descriptions
%    - New §16: Hint System (hints.js, ? key toggle, 11 rules)
%    - §15 keyboard shortcuts: added ? key
%    - Gauges: period display limited to 3 significant figures
%    - New §10.7: Pause and Timeline system (timeline.js)
%    - §10.5 Keyboard shortcuts: added ← → while paused
%    - §14.2 Right Wing: added btnPause and btnCollect (moved from expert)
%    - §15.3 System: removed btnCollect, added btnThemeStyle
%    - §16 Hint System: rule table extended with six expert-panel-aware
%             rules; aggregate_formed split into _no_kick / _kick variants
%    - Appendix B removed
%    - §6.1 / §7.1: formation spread checks now use live measured
%             population spread (all floating particles for aggregates,
%             all live aggregates for pebbles) rather than the v_t spread
%             selector value; Appendix A: aggregate.minSpread removed
%    - §10.7.2: Game/Challenge modes disable timeline branching;
%             historical dots rendered greyed out
%    - §11.5: drawVtDistribution updated — sigma/mu readout removed;
%             two spread bars added below x-axis (green = particles,
%             amber = aggregates) showing measured 2-sigma width with
%             formation threshold line
%    - §14.2 Right Wing: Info button split into Manual + Hints (two
%             separate buttons); btnTheme updated — long press cycles
%             visual theme (default → modern → pfeiffer)
%    - §15.3 System: btnThemeStyle removed (theme cycling moved to
%             long-press on right-wing btnTheme); btnSettings added
%             (settings overlay: nozzle focus + depth effect sliders)
%
%  TODO (icons):
%    - icons/btnSysAutoOmega.png may need updating if the SVG shape for
%      mode 2 (launch/amber) differs from the original two-state design.
%      Check getOmegaCtlIconSVG(0/1/2) in ui.js and re-rasterise if needed.
% =============================================================================

\documentclass[12pt, a4paper]{article}

\usepackage{amsmath, amssymb}
\usepackage{geometry}
\usepackage{graphicx}
\usepackage[colorlinks=true, linkcolor=blue, urlcolor=blue]{hyperref}
\usepackage{listings}
\usepackage{xcolor}
\usepackage{booktabs}
\usepackage{float}
\usepackage{parskip}
\usepackage{array}
\usepackage{longtable}
\usepackage{caption}
\usepackage{subcaption}
\usepackage{mathtools}
\usepackage[T1]{fontenc}

\geometry{margin=2.5cm}

\newcommand{\btnicon}[1]{%
  \IfFileExists{icons/#1.png}{%
    \raisebox{-0.3ex}{\includegraphics[height=1.2em]{icons/#1}}%
  }{%
\raisebox{-0.1ex}{{\tiny\texttt{[#1]}}}%
  }%
}

\lstset{
  basicstyle=\ttfamily\small,
  keywordstyle=\color{blue},
  commentstyle=\color{gray},
  breaklines=true,
  frame=single,
  backgroundcolor=\color{gray!8},
}

\newcommand{\vt}{v_{\mathrm{t}}}
\newcommand{\xc}{x_{\mathrm{c}}}
\newcommand{\Hcx}{H_{\mathrm{cx}}}
\newcommand{\Hr}{H_{\mathrm{r}}}
\newcommand{\omg}{\omega}

% =============================================================================
\begin{document}
% =============================================================================

\pagenumbering{roman}
\begin{titlepage}
  \centering
  \vspace*{3cm}
  {\Huge\bfseries Levitation Lab\\[0.4em]}
  {\Large Technical Documentation\\[2em]}
  {\large GT4Pebbles Project\\[0.5em]}
  {\large \today}
  \vfill
  {\small
    This document describes the full physics model, numerical implementation,
    analytical tooling, user interface, and game design of the Levitation Lab
    browser simulation. It is intended for developers, researchers, and outreach
    coordinators who need to understand, extend, or verify the simulation.
  }
\end{titlepage}

\tableofcontents
\clearpage
\pagenumbering{arabic}
\setcounter{page}{1}


% =============================================================================
\section{Introduction}
% =============================================================================

The \textbf{GT4Pebbles} (Ground Truth for Pebbles) project studies how planet
formation begins in protoplanetary disks around young stars. A central
experimental apparatus is a rotating drum in which dust particles are
levitated for extended periods, allowing them to collide and grow through a
sequence of stages: free particles $\to$ aggregates $\to$ compressed
aggregates---a prestage for pebble formation.

The \textbf{Levitation Lab} is a browser-based interactive simulation
that captures the core physics of this experiment. It is used as a
game, as an outreach tool---allowing members of the public to
explore the experiment, and as a pedagogical tool for students and
researchers. The simulation is built entirely in vanilla JavaScript
and HTML5 Canvas, with no external physics engine.

This document covers:
\begin{itemize}
  \item The underlying physical model and all governing equations.
  \item The numerical implementation, including timestep management and
        substep logic.
  \item Every user interface control, including button icons, with rationale
        for design choices.
  \item The game and challenge modes and their win/fail logic.
  \item The expert analytical panel and its visualisation tools.
  \item Audio, viewport, and URL parameter systems.
\end{itemize}

% =============================================================================
\section{Coordinate System and Drum Geometry}
% =============================================================================

\subsection{Coordinate System}

The simulation uses a right-handed coordinate system centred on the drum axis,
with $x$ pointing right and $y$ pointing upward. Positions are measured in
\textbf{drum units} (equivalent to centimetres in the physical experiment).

The drum has inner radius $R = 100\,\mathrm{du}$ and a hub of radius
$R_{\mathrm{axis}} = 5\,\mathrm{du}$. The drum extrusion depth (into the
screen) is $20\,\mathrm{du}$, but this dimension is not used in the 2-D
simulation; it serves only as a conceptual reference.

Canvas rendering converts drum coordinates to screen pixels as:
\begin{align}
  x_{\mathrm{px}} &= C_x + x \cdot S, \\
  y_{\mathrm{px}} &= C_y - y \cdot S,
\end{align}
where $(C_x, C_y)$ is the screen-space drum centre and $S$ is the pixel scale
(pixels per drum unit), both computed by the layout engine from the window size.

\subsection{Levitation Zone}

The levitation zone (the \emph{highlight zone} in the code) is a disc centred
at $(\Hcx, H_{\mathrm{cy}}) = (50, 0)$ with radius $\Hr = 50\,\mathrm{du}$.
This places the zone on the right side of the drum, tangent to both the drum
wall and the axis. A particle is considered \emph{levitated} when it has
remained continuously inside the levitation zone for at least one full drum
revolution.

% =============================================================================
\section{Drum Rotation Dynamics}
% =============================================================================

\subsection{Angular Velocity Model}

The drum has two angular velocity variables: the \emph{target} velocity
$\omg_{\mathrm{target}}$ (set by the user via swipe or auto-omega) and the
\emph{actual} velocity $\omg$.

The target decays exponentially each frame:
\begin{equation}
  \omg_{\mathrm{target}}(t + \Delta t) = \omg_{\mathrm{target}}(t)
    \cdot e^{-\lambda_d \Delta t},
\end{equation}
with $\lambda_d = 1/30\,\mathrm{s}^{-1}$ (i.e.\ a 30-second decay time
constant). This models aerodynamic drag on the drum once the user stops
spinning it. The decay can be disabled with the \textbf{$\Omega$ lock} button.

The actual velocity tracks the target with a slip rate
$\lambda_s = 0.7\,\mathrm{s}^{-1}$:
\begin{equation}
  \omg(t + \Delta t) = \omg(t) + \bigl[\omg_{\mathrm{target}}(t+\Delta t)
    - \omg(t)\bigr] \cdot \bigl(1 - e^{-\lambda_s \Delta t}\bigr).
\end{equation}

\subsection{User Input: Swipe}

A pointer drag across the drum surface imparts angular momentum proportional
to the tangential component of the drag velocity:
\begin{equation}
  \Delta\omg_{\mathrm{target}} = g_s \cdot
    \frac{\Delta s_{\mathrm{tang}}}{r / R},
\end{equation}
where $g_s = 0.003\,\mathrm{(rad\,s^{-1})/px}$ is the swipe gain,
$\Delta s_{\mathrm{tang}}$ is the pointer displacement projected onto the
tangential direction, and $r$ is the distance of the midpoint from the axis.
The denominator $(r/R)$ normalises for swipes near the axis versus near the
wall, giving consistent feel at all radii. The result is clamped to
$|\omg_{\mathrm{target}}| \leq \omg_{\max} = 2\pi \times 1.5\,\mathrm{rad\,s^{-1}}$.

\paragraph{Design note.} The physical experiment runs at approximately
$1\,\mathrm{RPM} \approx 0.105\,\mathrm{rad\,s^{-1}}$. At that speed,
the levitation feedback to the user would be far too slow for an interactive
experience. The simulation maximum of $1.5\,\mathrm{rev\,s^{-1}}$ is
approximately 90 times faster, but preserves all the qualitative physics.

% =============================================================================
\section{Particle Kinematics and Levitation}
% =============================================================================

\subsection{Equations of Motion}

In the rotating frame, the gas co-rotates with the drum and exerts a drag force
on the particle that brings its velocity to the local gas velocity minus its
terminal settling velocity $\vt$. In the (non-rotating) lab frame, the
particle's instantaneous velocity is:
\begin{align}
  \dot{x} &= -\omg y, \label{eq:vx} \\
  \dot{y} &= +\omg x - \vt. \label{eq:vy}
\end{align}

This is a \emph{kinematic} model: accelerations are not integrated; the
velocity is set to the equilibrium value at every substep. This is valid when
the aerodynamic drag time is much shorter than the orbital period, which is
the overdamped (high-Stokes-drag) limit appropriate for the small particles
of interest.

\subsection{Circular Orbits}

Setting $\ddot{x} = \ddot{y} = 0$ gives the orbit centre:
\begin{equation}
  \boxed{\xc = \frac{\vt}{\omg}, \qquad y_c = 0.}
  \label{eq:orbit_centre}
\end{equation}

A particle at $(x_0, y_0)$ orbits a circle of radius
$r = \sqrt{(x_0 - \xc)^2 + y_0^2}$
centred at $(\xc, 0)$ at angular rate $\omg$.

\subsection{Terminal Velocity Distributions}

Each injection event draws $N_P$ particles with terminal velocities from one
of three distributions:

\subsubsection{Gaussian (default)}
\begin{equation}
  \vt^{(i)} = \bar{\vt} \cdot \bigl(1 + s \cdot \xi_i\bigr),
  \qquad \xi_i \sim \mathcal{N}(0, 1),
\end{equation}
where $\bar{\vt}$ is the mean terminal velocity and $s$ is the relative spread
parameter ($0 \leq s \leq 0.5$).

\subsubsection{Bi-disperse}
Two groups with terminal velocities $\vt^{(1)}, \vt^{(2)}$, fractional
spreads $s_1, s_2$, and mixing ratio $q$:
$P(\text{group 1}) = q/(1 + q)$.

\subsubsection{Power-law}
$n(v) \propto v^q$ on $[v_{\min}, v_{\max}]$, sampled by inverse-transform.
Special case $q = -1$: log-uniform. All $\vt$ values are clamped from below
at the settle floor $= 2\,\mathrm{cm\,s^{-1}}$.

\subsection{Visual Size}

The rendered radius scales as $s_{\mathrm{vis}} = \sqrt{\vt / 30}$,
clamped to $[0.55, 1.8]\,\mathrm{du}$, and reduced by a factor $0.1$--$0.7$
for $N_P \geq 300$.

\subsection{Injection and Release}

Particles are released at $y = 110\,\mathrm{du}$, spread uniformly in $x$
across $[0, 100]\,\mathrm{du}$. They fall freely until crossing the drum wall,
at which point the kinematic drag model activates.

\subsection{Wall Collision and Particle Loss}

When $r \geq R - r_{\mathrm{coll}}$ ($r_{\mathrm{coll}} = 0.5\,\mathrm{du}$),
a particle is stuck to the wall, displayed in red, fades over $5\,\mathrm{s}$,
and removed. With \textbf{Nanocoating} active, it is elastically reflected
instead.

% =============================================================================
\section{Automatic Drum Speed Control}
% =============================================================================

\subsection{Orbit-Centring Formula}

Centring the orbit-centre range $[\vt^{\min}/\omg,\; \vt^{\max}/\omg]$
on $\Hcx = 50$ gives:
\begin{equation}
  \boxed{\omg_{\mathrm{auto}} = \frac{\vt^{\min} + \vt^{\max}}{2 \Hcx}.}
  \label{eq:auto_omega}
\end{equation}

\subsection{Launch Mode}

In \emph{launch mode}, the controller waits for injection, estimates the time
of peak particle flux using a KDE of arrival times, spins the drum up just
before that time, then transitions to live tracking.

% =============================================================================
\section{Aggregate Formation}
% =============================================================================

\subsection{Formation Threshold and Trigger}

Formation requires $\geq 30$ levitated particles and a measured
$v_t$ spread of the \emph{all floating} particle population of
$\geq 20\%$ (assessed as the 95.45th-percentile width of the live
distribution, independent of the $v_t$ spread selector value).
A revolution counter accumulates at $\Delta t / T$; the first aggregate forms
after 3 revolutions, subsequent ones after 1 revolution.

\subsection{Merge Animation}

10 levitated particles near a randomly chosen centre are animated toward
their centroid over $0.6\,\mathrm{s}$, then replaced by a single aggregate.

\subsection{Aggregate Terminal Velocity and Stokes Number}

The aggregate's terminal velocity is the mass-weighted mean of its constituent
monomers. When the \textbf{Stokes kick} is enabled, a factor
$f_{\mathrm{Stokes}} = 1.2$ is applied:
\begin{equation}
  \vt^{\mathrm{agg}} = \bar{\vt} \cdot f_{\mathrm{Stokes}}.
\end{equation}
This reflects two physical effects: increased inertia (higher Stokes number)
and cross-section shadowing (reduced drag per unit mass), both of which shift
the aggregate's orbit centre $\xc = \vt/\omg$ slightly outward relative to
the monomer population. When Stokes kick is off ($f_{\mathrm{Stokes}} = 1$),
the aggregate orbits at exactly the mass-weighted mean monomer orbit centre.

\subsection{Aggregate Orbital Mechanics and Brownian Perturbations}

Aggregates follow the same kinematic orbit equations as monomers.
In growth mode, Langevin perturbations ($\sigma_B = 2\,\mathrm{cm\,s^{-1}}$,
restoring spring $k_R = 0.5\,\mathrm{s}^{-1}$) cause orbital drift that
enables close approaches.

\subsection{Aggregate--Aggregate Growth Collisions}

In growth mode, overlapping pairs merge over $0.6\,\mathrm{s}$. An aggregate
reaching $n \geq 100$ monomers becomes a pebble.

\subsection{Wall Collision of Aggregates}

Wall contact fragments the aggregate into 10 stuck monomers. With Nanocoating,
it is elastically reflected instead.

% =============================================================================
\section{Pebble (Golden Ball) Formation}
% =============================================================================

\subsection{Formation Threshold}

7 levitated aggregates held for 3 revolutions, with a measured $v_t$ spread
of \emph{all live aggregates} $\geq 30\%$ (assessed as the
95.45th-percentile width of the live aggregate distribution, independent of
the $v_t$ spread selector value).
The densest cluster is found by K-nearest-neighbours ($K=3$).

\subsection{Golden Ball Physics}

Newtonian dynamics under $g = 200\,\mathrm{du\,s^{-2}}$, wall restitution
$e_w = 0.85$, friction $\mu_w = 0.92$, ball--ball restitution $e_b = 0.90$.
A rotating bump on the drum wall adds energy. Wall-contact drag interpolates
toward drum surface velocity with $\lambda_c = 0.7\,\mathrm{s}^{-1}$.

% =============================================================================
\section{Globe (Planet) Formation}
% =============================================================================

3 pebbles merge into a globe over $1.2\,\mathrm{s}$; up to 7 simultaneous.
Globes are attracted toward $(0, -50)$ by a spring, repelled from each other
by a $1/r^2$ force ($F_{\mathrm{rep}} = 15000\,\mathrm{du^3\,s^{-2}}$), and
damped at $\exp(-9.74\,\Delta t)$ per frame.

% =============================================================================
\section{Solar System Phase Machine}
% =============================================================================

On the second globe: \texttt{showing} $\to$ \texttt{transitioning}
$\to$ \texttt{orbiting} $\to$ \texttt{spindown} $\to$ \texttt{final\_move}
$\to$ \texttt{final\_view}. Modified Keplerian law
$\omg_{\mathrm{orbit}} = \omg_0 / r^{0.75}$ with inclination $\cos i = 0.3$.
Two Voyager probes are launched after the system stabilises.

% =============================================================================
\section{Numerical Implementation}
% =============================================================================

\subsection{Timestep and Substep Logic}

Frame duration capped at $\Delta t_{\max} = 0.033\,\mathrm{s}$, then scaled
by $s_{\mathrm{sim}}$. Physics advances in substeps:
\begin{equation}
  n_{\mathrm{sub}} = \max\Bigl(1,\; \left\lceil\frac{\Delta t}{\Delta t_{\mathrm{sub}}}\right\rceil\Bigr),
  \qquad \Delta t_{\mathrm{sub}} = \frac{2\pi/\omg_{\max}}{100}.
\end{equation}

\subsection{Slow-Motion During Merges}

When slow-motion is armed, any frame with an active merge advances at
$\Delta t / 5$.

\subsection{Stroboscopic Mode}

Canvas redraws once per drum revolution; the revolution counter is
$\lfloor |\theta_{\mathrm{drum}}| / 2\pi \rfloor$.

\subsection{Simulation Speed Gears}

Ten speed gears are available:
$\{1/64\times,\, 1/32\times,\, 1/16\times,\, 1/8\times,\, 1/4\times,\,
1/2\times,\, 1\times,\, 2\times,\, 4\times,\, 8\times\}$.
The display label uses tightened letter-spacing for sub-unity speeds
(e.g.\ \texttt{/64}) to fit the constrained button width.
The centre button shows the current gear and snaps to $1\times$ when pressed.
The default index is $1\times$; all physics including drum decay and
injection timing scale uniformly with the multiplier.

\subsection{Period Display}

The drum rotation period is displayed as $T = 2\pi/|\omg|$ (or $\infty$ when
$|\omg| < 10^{-3}\,\mathrm{rad\,s^{-1}}$), formatted to 3 significant figures
using \texttt{parseFloat(T.toPrecision(3))}, which strips trailing zeros.

\subsection{Pause}

\texttt{Space} (or the \textbf{Pause} button in the instrument block)
freezes all physics, drum rotation, and motor audio while leaving the canvas
rendering active. When paused, the timeline scrubber is revealed (see
\S\ref{sec:timeline}).

\subsection{Pause and Timeline System}
\label{sec:timeline}

The timeline system (\texttt{timeline.js}) maintains a rolling buffer of full
simulation snapshots and allows the user to branch the simulation from any
saved point.

\subsubsection{Snapshot Recording}

One snapshot is recorded every $\Delta t = 10\,\mathrm{s}$ of wall-clock time
while the simulation is running and not paused. Up to $N = 20$ snapshots are
retained, giving up to $200\,\mathrm{s}$ of history under normal conditions.

When the buffer exceeds either the count cap or a memory budget of
$20\,\mathrm{MB}$ (measured as a JSON proxy), snapshots are thinned by
repeatedly removing interior elements starting from index 1 and advancing
toward the most recent end. This produces logarithmic spacing: the gap between
adjacent snapshots at age $\tau$ before now grows as $\delta t(\tau) \propto
\tau$, so the number of snapshots covering a history span $T$ scales as
$\log(T/\Delta t)$. The oldest and most recent snapshots are always preserved;
only the middle is thinned. The practical effect is that recent history remains
at $10\,\mathrm{s}$ granularity while distant history degrades gracefully to
coarser intervals, rather than the oldest history being discarded entirely.

Each snapshot is a shallow copy of the complete simulation state:
\texttt{particles}, \texttt{aggregates}, \texttt{goldenBalls},
\texttt{globes}, \texttt{toInject}, scalar fields (\texttt{omega},
\texttt{drumAngle}, \texttt{t}, \texttt{lostCount}, \texttt{aggCount},
\texttt{eggBallCount}, hold-revolution counters), \texttt{solar}, and
\texttt{tray}. Particle trails and diagnostic-target flags are excluded.
Heatmap accumulators and process button states (\textsc{Nanocoating},
\textsc{Growth mode}, \textsc{Stokes kick}, etc.) are intentionally omitted.
On restore, the simulation resumes with whatever process configuration is
currently active --- branching from a past snapshot is intended as a way to
correct a mistake in process settings, not to replay the original trajectory
faithfully.

\subsubsection{Pause and Scrubber Display}

On pause, a \emph{frozen snapshot} of the current state is captured and stored
separately from the rolling buffer. A timeline bar is drawn on
\texttt{ctxOv} below the drum interior showing one dot per rolling snapshot
plus a final dot for ``NOW'' (the frozen snapshot). The current cursor
position is highlighted with a larger amber dot and a pulsing halo. A time
label below the cursor shows the offset from NOW (e.g.\ $-30\,\mathrm{s}$).

In \textbf{Game} and \textbf{Challenge} modes, timeline branching is
disabled. Historical snapshot dots are rendered greyed out, and navigating to
them is a no-op. This prevents competitive runs from being replayed or
rewound for score advantage.

\subsubsection{Navigation}

While paused:
\begin{itemize}
  \item \textbf{Keyboard}: $\leftarrow$ / $\rightarrow$ step the cursor one
    snapshot at a time.
  \item \textbf{Pointer}: tapping or clicking any dot (22\,px hit radius)
    jumps to that snapshot directly.
\end{itemize}
At each step the full simulation state is immediately restored from the
corresponding snapshot, updating the canvas.

\subsubsection{Branching on Resume}

When the user resumes (via \texttt{Space} or the Pause button):
\begin{itemize}
  \item If the cursor is at NOW, the simulation resumes from the frozen state
    with no history change.
  \item If the cursor is in the past, the rolling buffer is truncated to
    include only snapshots up to and including the cursor position. All
    snapshot \texttt{wallT} timestamps are then shifted forward by
    $(t_{\mathrm{now}} - t_{\mathrm{cursor}})$ so that the branch point reads
    as ``just now'' and older snapshots retain their correct relative spacing.
    \texttt{lastSnapWallT} is set to the current wall time. The simulation then
    runs forward from the branched state, overwriting history from that point.
\end{itemize}

\subsubsection{Merge State on Restore}

Any in-progress merge animation (\texttt{aggMerging}, \texttt{eggMerging},
\texttt{globeMerging}, \texttt{aggGrowMerging}) is set to \texttt{null} on
restore, and all objects have their \texttt{merging} flag cleared. Since the
conditions that triggered the original merge are preserved in the restored
state (levitated count, spread, hold counters), the merge will re-trigger
naturally within the next few revolutions.

% =============================================================================
\section{Heatmap and Analytical Overlays}
% =============================================================================

\subsection{Heatmap Accumulation}

$50 \times 50$ grid; per-cell accumulation of $n$, $\Sigma v$, $\Sigma v^2$
finalised once per revolution. Three modes: velocity dispersion $\sigma_v$,
particle density $\rho$, collision proxy $\mathcal{C} = \rho \sigma_v^2$.

\subsection{Orbit Projections}

Selected orbits mode: orbit circles centred at $(\vt/\omg, 0)$, sample
updated every $3\,\mathrm{s}$.

\subsection{Aggregate Encounter Prediction}

For each pair of live, free-floating aggregates, the closest-approach position
during the next orbit is computed analytically. The squared separation as a
function of phase advance $\theta$ is:
\begin{equation}
  d^2(\theta) = C_0 + 2\,\Delta\xc\,(P\cos\theta - Q\sin\theta),
\end{equation}
minimised at $\theta^* = \arctan(-Q, P)$ (both candidates evaluated).

The number of pairs tracked is \texttt{TUNING.encounters.showPairs} (default
10). Colour encodes collision prediction:
\begin{itemize}
  \item \textbf{Green} --- $d < r_i + r_j$: collision predicted this orbit.
    Takes precedence at any rank.
  \item \textbf{Amber} --- top \texttt{arcPairs} pairs (default 2) that are
    not colliding.
  \item \textbf{Gray-purple} --- all remaining tracked pairs.
\end{itemize}
When a pair is within $\Delta\theta < 0.15\,\mathrm{rad}$ of closest approach,
the indicator brightens subtly in its own hue. For the top \texttt{arcPairs}
pairs, an arc is drawn on each aggregate's orbit from its current position to
the predicted closest-approach point in the direction of motion. Cache updates
at $\approx 5\,\mathrm{Hz}$; arc directions recomputed every render frame.

\subsection{Ghost Paths}

Selected particle: white solid (actual orbit), cyan dashed (gas co-rotation),
purple dashed (free-fall vacuum). Click any particle to select it.

\subsection{\texorpdfstring{$v_t$}{vt} Distribution Display}

The distribution panel shows three overlapping visualisations drawn on the
same axes:
\begin{itemize}
  \item \textbf{KDE curves} --- yellow (all floating particles), green
    (levitated particles), amber histogram (aggregates).
    Silverman bandwidth $h = 1.06\,\hat\sigma\,n^{-1/5}$. All normalised
    independently.
  \item \textbf{Particle spread bar} (green, below the $x$-axis) --- shows
    the measured $2\sigma$ width of the floating particle population as a
    horizontal bar, with a vertical threshold line at the $20\%$ formation
    threshold. The bar gives an immediate visual indication of whether the
    current population is wide enough to form aggregates.
  \item \textbf{Aggregate spread bar} (amber, below the particle bar) ---
    shows the measured $2\sigma$ width of the live aggregate population,
    with a threshold line at $30\%$. Indicates proximity to pebble formation.
\end{itemize}
Both spread bars use the same 95.45th-percentile range estimator that governs
actual formation checks (see \S6.1 and \S7.1), so the bar reaching the
threshold line is the exact condition that gates formation.

\subsection{Aggregate Size Histogram}

Bar chart of monomer count distribution (bins 10--100), plus pebble count.
$y$-axis fixed at 10; count labels above bars.

\subsection{Live KDE Orbit Band}

The orbit centre projection button (\btnicon{btnVtProj}) has three states:

\begin{enumerate}
  \item \textbf{Off} (default dark cap) --- nothing drawn.
  \item \textbf{Predicted} (amber) --- a dashed line on the drum equator
    spanning $[\vt^{\min}/\omg,\; \vt^{\max}/\omg]$ with dots at the
    endpoints and centre, computed from the current parameter selectors.
    This is a pure prediction and does not depend on particles being present.
  \item \textbf{Live} (green) --- a KDE-shaded band, 4 drum-units tall,
    centred on the equator and clipped to the drum interior. Cyan represents
    all free-floating particles; amber represents free-floating aggregates.
    The alpha at each horizontal position is proportional to the Gaussian KDE
    of the respective population's orbit centres
    $\xc^{(i)} = \vt^{(i)} / \omg$.
    When aggregates are present, the band splits: particles occupy the lower
    half, aggregates the upper half.
    Bandwidth uses Silverman's rule:
    $h = \max(0.5,\; 1.06\,\hat\sigma\,n^{-0.2})$ drum units.
\end{enumerate}

% =============================================================================
\section{Collection Tray}
% =============================================================================

\subsection{Arming and Insertion}

Blade extends from the slot marker inward, sweeping $240^{\circ}$ ($4\pi/3$) in two
phases: cruise ($2/3$, constant speed) and brake ($1/3$, cubic smoothstep to
zero). Brake duration computed dynamically from $\omg_0$.
(\texttt{TUNING.tray.brakeDur} is vestigial and not used at runtime.)

\subsection{Levitation State Invariant on Capture}

When any object is pinned to the tray by \texttt{tryCatchOnTray}, its
\texttt{inHighlightSince} timestamp is immediately set to \texttt{null},
preventing stale levitation timestamps. The per-frame particle update in
\texttt{step()} also clears \texttt{inHighlightSince} for any particle
already \texttt{onTray}. All levitation-querying functions
(\texttt{eggLevitatedParticles}, \texttt{eggLevitatedAggregates},
\texttt{resolveAggAggCollisions}, \texttt{\_updateEncounterCache})
explicitly exclude \texttt{onTray} objects.

% =============================================================================
\section{Slice Widening (Duplication)}
% =============================================================================

Pressing \btnicon{btnProcDouble} duplicates all levitated particles and
aggregates, placing each copy $90^{\circ}$ ahead in its orbit:
\begin{align}
  x' &= \xc - y, \quad y' = x - \xc, \\
  v_x' &= -\omg y', \quad v_y' = +\omg x' - \vt.
\end{align}
Copies inherit \texttt{inHighlightSince} from the originals.

% =============================================================================
\section{User Interface --- Main Controls}
% =============================================================================

\subsection{Left Wing: Injection Parameters}

\begin{longtable}{@{}p{2cm}lp{9cm}@{}}
  \toprule
  Icon & Label & Description \\
  \midrule
  \btnicon{winNP}  & \textbf{Particles} &
    $N_P \in \{1,2,3,10,30,100,300,1000,3000,10000\}$.\\
  \btnicon{winVT}  & \textbf{$v_t$ cm/s} &
    $\bar{\vt} \in \{2,5,10,20,30,40,50\}\,\mathrm{cm\,s^{-1}}$.\\
  \btnicon{winSpread} & \textbf{$v_t$ spread} &
    $s \in \{0,0.01,...,0.50\}$. A measured population spread of at
    least 20\% is needed for aggregates; 30\% for pebbles.\\
  \btnicon{winDT}  & \textbf{$\Delta t$ inject} &
    $\Delta t_{\mathrm{inj}} \in \{0,0.5,...,10\}\,\mathrm{s}$.\\
  \bottomrule
\end{longtable}

Clicking the VT or Spread display when a custom distribution profile is active
opens the distribution overlay directly.

\subsection{Right Wing: Action Controls}

\begin{longtable}{@{}p{2cm}lp{9cm}@{}}
  \toprule
  Icon & Label & Description \\
  \midrule
  \btnicon{btnStart}    & \textbf{Inject}      & Release a new batch. \\
  \btnicon{btnSound}    & \textbf{Sound}       & Toggle audio. \\
  \btnicon{btnOmegaCtl} & \textbf{$\Omega$}    & Lock drum rotation rate. \\
  \btnicon{btnTrails}   & \textbf{Trails}      & Particle trail rendering ($N_P < 300$). \\
  \btnicon{btnLaser}    & \textbf{Lidar}       &
    Lidar mode; laser at $\omg_{\mathrm{laser}} = 2\pi/1.2\,\mathrm{rad\,s^{-1}}$. \\
  \btnicon{btnMicro}    & \textbf{Microscope}  & Slide viewport down --- microscope channel. \\
  \btnicon{btnChart}    & \textbf{Chart}       & Slide viewport down --- chart channel. \\
  \btnicon{btnPause}    & \textbf{Pause}       &
    Freeze simulation; reveal timeline scrubber (see \S\ref{sec:timeline}).
    Also toggled by \texttt{Space}. \\
  \btnicon{btnCollect}  & \textbf{Collect}     &
    Arm collection tray (see \S12). Pulses amber while inserting.
    Automatically disables auto-omega when the brake phase begins,
    preventing the controller from fighting the tray deceleration. \\
  \btnicon{btnTheme}    & \textbf{Lighting}    &
    Short press: toggle dark / light lab background.
    Long press: cycle visual theme (default brass-and-dark $\to$
    clinical \emph{modern} $\to$ \emph{Pfeiffer} red-structures variant
    $\to$ default). \\
  \btnicon{btnGameMode} & \textbf{Game}        & 8-level progressive game mode. \\
  \btnicon{btnChallenge}& \textbf{Challenge}   & Timed challenge with leaderboard. \\
  \btnicon{btnManual}   & \textbf{Manual}      & Open user manual. \\
  \btnicon{btnHints}    & \textbf{Hints}       &
    Toggle the hint system (see \S\ref{sec:hints}). Equivalent to pressing
    \texttt{?}. \\
  \btnicon{btnReset}    & \textbf{Reset}       &
    Clear all objects; reset to clean state. Also disables zoom mode if active. \\
  \bottomrule
\end{longtable}

% =============================================================================
\section{Expert Panel}
% =============================================================================

Opened by triple-tapping the title plate or appending \texttt{?expert}.

\subsection{HUD --- Visualisation Overlays}

\begin{longtable}{@{}p{1.5cm}lp{9.5cm}@{}}
  \toprule
  Icon & Name & Description \\
  \midrule
  \btnicon{btnVtProj} & \textbf{Orbit centre / KDE band} &
    Three-state button.
    \emph{Amber} (state 1): predicted dashed line on the equator spanning
    $[\vt^{\min}/\omg,\; \vt^{\max}/\omg]$ from the parameter selectors.
    \emph{Green} (state 2): live KDE band --- cyan for particles, amber for
    aggregates, alpha proportional to local orbit-centre density. When
    aggregates are present the band splits: particles below, aggregates above
    the equator. Off by default. \\
  \btnicon{btnFlashOrbit} & \textbf{Flash orbit} &
    After each aggregate forms, briefly draws its orbit ring alongside the
    mean-$\vt$ levitated orbit, connected by a line showing any Stokes kick
    offset. \\
  \btnicon{btnAggEncounters} & \textbf{Encounter projections} &
    Closest-approach positions for the top
    \texttt{TUNING.encounters.showPairs} pairs (default 10). Green if
    collision predicted ($d < r_i + r_j$), amber for the top
    \texttt{arcPairs} non-colliding pairs, gray-purple for the rest.
    Green takes precedence at any rank. Top \texttt{arcPairs} pairs also
    show orbit arcs from current position to closest approach. \\
  \btnicon{btnVtDist} & \textbf{$v_t$ distribution} &
    Live KDE with spread bars: levitated (green), floating (yellow),
    aggregate histogram (amber), particle spread bar (green, below axis),
    aggregate spread bar (amber, below axis). \\
  \btnicon{btnAggHist} & \textbf{Size histogram} &
    Aggregate monomer count distribution + pebble bar. \\
  \btnicon{btnSysSolidMap} & \textbf{Backdrop} &
    Opaque backing for contrast. Long press opens opacity slider. \\
  $\blacktriangleleft$ & \textbf{Mode prev} & Previous analysis mode. \\
  \btnicon{btnHudMaster} & \textbf{Analysis HUD} &
    Master switch; six modes: vector field, ghost paths, orbit sample,
    density, dispersion, collision proxy. \\
  $\blacktriangleright$ & \textbf{Mode next} & Next analysis mode. \\
  \bottomrule
\end{longtable}

\subsection{Processes}

\begin{longtable}{@{}p{1.5cm}lp{9.5cm}@{}}
  \toprule
  Icon & Name & Description \\
  \midrule
  \btnicon{btnProcInvinc} & \textbf{Nanocoating} & No wall losses. \\
  \btnicon{btnProcPeb}    & \textbf{Sweep}        & Remove non-levitated particles instantly. \\
  \btnicon{btnProcDouble} & \textbf{Duplicate}    & Duplicate levitated objects $90^{\circ}$ ahead. \\
  \btnicon{btnProcPureV}  & \textbf{Stokes kick}  & $\times 1.2$ $\vt$ at formation, $\times 1.05$ per growth merge. \\
  \btnicon{btnProcAgg}    & \textbf{Aggregate formation} & Green = on; red = disabled. \\
  \btnicon{btnProcGrowth} & \textbf{Growth mode}  &
    Green: pebble path. Amber: agg--agg collision growth. Red: off. \\
  \bottomrule
\end{longtable}

\subsection{System}

\begin{longtable}{@{}p{1.5cm}lp{9.5cm}@{}}
  \toprule
  Icon & Name & Description \\
  \midrule
  \btnicon{btnDistMenu}      & \textbf{Distribution}    & Gaussian / bi-disperse / power-law profile. \\
  \btnicon{btnSysAutoOmega}  & \textbf{Cruise / Autopilot} &
    Three-state button. \emph{Off}: manual control.
    \emph{Green (auto)}: live orbit-centring via Eq.~\ref{eq:auto_omega}.
    \emph{Amber (launch)}: waits for injection, spins up at peak flux, then
    transitions to auto tracking. \\
  \btnicon{btnSettings}      & \textbf{Settings}        &
    Opens the settings overlay. Contains a \emph{nozzle focus} slider
    (controls the lateral spread of the injection nozzles) and a
    \emph{depth effect} slider (controls the intensity of the cylindrical
    3-D rendering on the drum band; $0$ = flat, $1$ = full depth). \\
  $\blacktriangleleft$       & \textbf{Speed down}      & Reduce simulation speed one gear. \\
  \btnicon{btnSpeedReset}    & \textbf{Speed}           &
    Current gear label; press to reset to $1\times$. \\
  $\blacktriangleright$      & \textbf{Speed up}        & Increase simulation speed one gear. \\
  \btnicon{btnSysStrobe}     & \textbf{Strobe}          & Redraw once per drum revolution. \\
  \btnicon{btnProcSlowMo}    & \textbf{Slow motion}     &
    Short press: $5\times$ slow during merges.
    Long press ($>2\,\mathrm{s}$): Fast Chain mode --- all formation thresholds
    lowered for rapid testing. \\
  \btnicon{btnProcFast}      & \textbf{Zoom}            & Levitation zone fills drum area. \\
  \bottomrule
\end{longtable}

\subsection{Keyboard Shortcuts}

\begin{longtable}{@{}lp{11cm}@{}}
  \toprule
  Key & Effect \\
  \midrule
  \texttt{Space} &
    Pause and resume the simulation. Drum rotation, all physics, and motor
    audio freeze; the timeline scrubber appears (see \S\ref{sec:timeline}). \\
  $\leftarrow$ / $\rightarrow$ &
    While paused: step the timeline cursor one snapshot backward or forward.
    The canvas updates immediately. On resume from a past snapshot the future
    history is discarded and the simulation branches forward. \\
  \texttt{?} &
    Toggle the hint system (see \S\ref{sec:hints}). When on, contextual
    observations appear inside the drum. Suppressed during Game and Challenge
    modes and during merge animations. \\
  \bottomrule
\end{longtable}

% =============================================================================
\section{Hint System}
\label{sec:hints}
% =============================================================================

The hint system is a rule-based observer that displays brief contextual
messages inside the drum when the current simulation state warrants commentary.
It is implemented in \texttt{hints.js} and is independent of the game and
challenge modes, which suppress it entirely.

\subsection{Toggle and Display}

Pressing \texttt{?} (or the \textbf{Hints} button in the right wing) toggles
\texttt{window.hintsOn}. Messages appear as italic brass-coloured text on a
dark rounded backdrop near the top of the drum interior
(drum coordinate $y \approx 78$), drawn on \texttt{ctxOv} each frame
independently of strobe mode. Long messages are word-wrapped onto two lines.
The opacity follows a three-phase envelope:
fade-in $0.5\,\mathrm{s}$, hold $5.5\,\mathrm{s}$, fade-out $1.0\,\mathrm{s}$.

A global cooldown of $12\,\mathrm{s}$ prevents consecutive hints from
overlapping. No hint fires during an active merge animation.

\subsection{Rule Structure}

Each rule is a plain object:
\begin{lstlisting}[language=Java]
{
  id:       'unique_id',
  priority: 65,          // higher fires first
  cooldown: 45,          // seconds before this rule can fire again
  when: ctx => ...,      // condition on the context object
  message:  'Text.',     // displayed string
}
\end{lstlisting}

Rules that reference expert-panel buttons are gated on the helper
\texttt{\_expertOpen()}, which checks whether \texttt{expertContainer} carries
the \texttt{open} CSS class. These rules only fire when the expert door is
open, ensuring hints about expert controls are never shown to users who have
not discovered the panel.

The context object \texttt{ctx} passed to each \texttt{when} function exposes:
\texttt{particles.\{floating, levitated, lost, total, pending\}},
\texttt{aggregates.\{active, levitated, count\}},
\texttt{pebbles.\{count\}},
\texttt{drum.\{spinning, omega, period\}},
\texttt{params.\{VT\_SPREAD, N\_P, V\_T\}},
\texttt{merging}, \texttt{running}, \texttt{state\_t}.

\subsection{Rule Set}

Rules are listed in descending priority order. The ``Expert only'' column
marks rules gated on \texttt{\_expertOpen()}.

\begin{longtable}{@{}p{3.8cm}rcp{5.2cm}@{}}
  \toprule
  Rule id & Pri. & Expert & Message \\
  \midrule
  \texttt{pebble\_formed}              & 100 & --- &
    ``A pebble has formed --- seven aggregates compressed into one.'' \\
  \texttt{aggregate\_formed\_no\_kick} & 90  & --- &
    ``Aggregate formed. It orbits at the mean $v_t$ of its constituents.'' \\
  \texttt{aggregate\_formed\_kick}     & 90  & --- &
    ``Aggregate formed. Its higher Stokes number shifts its orbit outward.'' \\
  \texttt{pebble\_imminent}            & 85  & --- &
    ``Seven aggregates levitated. A pebble should form soon.'' \\
  \texttt{agg\_size\_dist}             & 82  & \checkmark &
    ``A spread of aggregate sizes has formed. Try the size histogram in the HUD.'' \\
  \texttt{suggest\_slow\_mo}           & 76  & \checkmark &
    ``A merge is building --- enable slow motion to watch it form.'' \\
  \texttt{spread\_low\_for\_pebble}    & 75  & --- &
    ``Pebble formation needs a $v_t$ spread of at least 30\%.'' \\
  \texttt{suggest\_encounters}         & 63  & \checkmark &
    ``Aggregates drifting together --- try the encounter projection HUD.'' \\
  \texttt{orbits\_crossing}            & 65  & --- &
    ``Orbits are crossing. Keep particles levitated and aggregates will form.'' \\
  \texttt{suggest\_vt\_dist}           & 53  & \checkmark &
    ``Particles sorting by size --- try the $v_t$ distribution HUD.'' \\
  \texttt{spread\_low\_for\_agg}       & 60  & --- &
    ``Aggregate formation needs a $v_t$ spread of at least 20\%.'' \\
  \texttt{monodisperse}                & 55  & --- &
    ``Same $v_t$ for all --- orbits never cross. Increase spread for collisions.'' \\
  \texttt{suggest\_orbit\_proj}        & 47  & \checkmark &
    ``Try the orbit projection button to see each size's orbit.'' \\
  \texttt{suggest\_auto\_omega}        & 44  & \checkmark &
    ``Auto-omega can centre the orbits automatically.'' \\
  \texttt{too\_many\_losses}           & 50  & --- &
    ``Too many wall losses. Try adjusting drum speed to centre the orbits.'' \\
  \texttt{spinning\_no\_particles}     & 40  & --- &
    ``Drum is spinning. Press Inject to release particles.'' \\
  \texttt{particles\_not\_spinning}    & 35  & --- &
    ``Particles are falling. Swipe the drum to spin it.'' \\
  \texttt{expert\_panel}              & 10  & --- &
    ``Expert controls are behind the locked door, bottom left.'' \\
  \bottomrule
\end{longtable}

The two aggregate-formation rules (\texttt{aggregate\_formed\_no\_kick} and
\texttt{aggregate\_formed\_kick}) share priority 90 but are mutually
exclusive by their \texttt{window.stokesKickOn} condition, so only one fires
at a time. This ensures the message is physically accurate.

% =============================================================================
\section{Game Mode}
\label{sec:game}
% =============================================================================

\begin{longtable}{@{}rp{5cm}p{9cm}@{}}
  \toprule
  \# & Name & Parameters / Goal / Fail \\
  \midrule
  1 & First Contact &
    $N_P=3$, $\vt=20$, $s=0$, $\Delta t=1\,\mathrm{s}$ \newline
    \textit{Goal:} $\geq 1$ levitated for $3\,\mathrm{s}$ \\[4pt]
  2 & Finding the Sweet Spot &
    $N_P=3$, $\vt=30$, $s=0$, $\Delta t=1\,\mathrm{s}$, trails \newline
    \textit{Goal:} $\geq 2$ levitated for $5\,\mathrm{s}$ \\[4pt]
  3 & Reading the Orbits &
    $N_P=5$, $\vt=50$, $s=0$, $\Delta t=1\,\mathrm{s}$, trails \newline
    \textit{Goal:} $\geq 2$ levitated for $5\,\mathrm{s}$ \\[4pt]
  4 & A Spread of Sizes &
    $N_P=15$, $\vt=30$, $s=0.15$, $\Delta t=2\,\mathrm{s}$, trails \newline
    \textit{Goal:} $\geq 4$ levitated for $5\,\mathrm{s}$ \\[4pt]
  5 & Flying Blind &
    $N_P=20$, $\vt=30$, $s=0.15$, $\Delta t=2\,\mathrm{s}$, lidar \newline
    \textit{Goal:} $\geq 6$ levitated for $8\,\mathrm{s}$ \\[4pt]
  6 & Crowded Skies &
    $N_P=50$, $\vt=30$, $s=0.15$, $\Delta t=3\,\mathrm{s}$ \newline
    \textit{Goal:} $\geq 10$ levitated for $10\,\mathrm{s}$ \\[4pt]
  7 & Encouraging Collisions &
    $N_P=100$, $\vt=30$, $s=0.20$, $\Delta t=3\,\mathrm{s}$ \newline
    \textit{Goal:} $\geq 1$ aggregate; timeout 40 revolutions \\[4pt]
  8 & More Energetic Encounters &
    $N_P=300$, $\vt=30$, $s=0.30$, $\Delta t=3\,\mathrm{s}$, trails \newline
    \textit{Goal:} $\geq 1$ pebble; timeout 100 revolutions \\
  \bottomrule
\end{longtable}

Structure-forming levels (7--8): timeout $= n_{\mathrm{revs}} \times 3\,\mathrm{s}$.
Win confirmation delay: $5\,\mathrm{s}$.
Level 8 completion opens the free-play sandbox with level-8 parameters.

% =============================================================================
\section{Challenge Mode}
\label{sec:challenge}
% =============================================================================

\begin{equation}
  \mathrm{Score} = N_{\mathrm{part}} + 10\,N_{\mathrm{agg}}
    + 100\,N_{\mathrm{pebble}} + 1000\,N_{\mathrm{planet}}.
\end{equation}

A particle counts if levitated for $\geq 3$ full revolutions. Time limit
$= 60 + \Delta t_{\mathrm{inj}}$ seconds. Scores stored in
\texttt{localStorage} as \texttt{levitation\_highscores}; top 7 displayed.

% =============================================================================
\section{Viewport System}
\label{sec:viewport}
% =============================================================================

Brass-framed screen, slides in on Microscope/Chart press. Rule engine
evaluates at $5\,\mathrm{Hz}$; crossfade $0.4\,\mathrm{s}$, debounce
$1.0\,\mathrm{s}$. Live renderers registered via
\texttt{registerViewportRenderer(name, fn)}.

% =============================================================================
\section{Audio System}
% =============================================================================

\begin{longtable}{@{}lp{11cm}@{}}
  \toprule
  Sound & Description \\
  \midrule
  Motor hum      & Two detuned sawtooth oscillators at 60/60.6 Hz; pitch and gain track $|\omg|$. \\
  Tink           & $220 \to 90\,\mathrm{Hz}$ sine; wall collision, $N_P \leq 30$; 500 ms cooldown. \\
  Snap           & 9 accelerating sawtooth clicks at $2500$--$4100\,\mathrm{Hz}$; aggregate formation. \\
  Aggregate merge & Sine thud $600 \to 220\,\mathrm{Hz}$, pitch tracks aggregate size; growth merges. \\
  Crunch         & 500 random sawtooth bursts; pebble formation. \\
  Chime          & 880/1320 Hz two-note arpeggio; particle levitation; 150 ms cooldown. \\
  Golden chime   & C5-E5-G5-D6 arpeggio; gauge first appearance. \\
  Golden thud    & $160 \to 70\,\mathrm{Hz}$; ball wall/collision impacts. \\
  Tray beep      & 880 Hz square wave, $\approx 350\,\mathrm{ms}$, every $1.0\,\mathrm{s}$
                   while tray is armed or inserting. \\
  Mill start     & Five-component sequence (blade, hiss, grind, thud, body); Inject press. \\
  Expert door    & Heavy clunk + 3-second cavernous reverb tail; door clicked while closed. \\
  \bottomrule
\end{longtable}

\texttt{AudioContext} resumed on first user gesture. Muting ramps gain to 0
over $50\,\mathrm{ms}$. Tab hide/show mutes and unmutes smoothly.

% =============================================================================
\section{URL Parameters}
\label{sec:url}
% =============================================================================

\begin{longtable}{@{}lp{11cm}@{}}
  \toprule
  Parameter & Effect \\
  \midrule
  \texttt{?expert}              & Open expert panel, activate $\Omega$ lock. \\
  \texttt{?game[=N]}            & Game Mode at level $N$ (1-indexed). \\
  \texttt{?challenge}           & Challenge Mode with default parameters. \\
  \texttt{?challenge\&np=N}     & Override challenge particle count (also \texttt{vt}, \texttt{spread}, \texttt{dt}). \\
  \texttt{?challenge\&time=N}   & Override challenge run length to $N$ seconds. \\
  \texttt{?modern}              & Clinical lab theme. \\
  \texttt{?pfeiffer}            & Pfeiffer variant (red structures). \\
  \texttt{?verify}              & Developer smoke-test: $N_P=3000$, $s=0.30$, $\Omega$ lock. \\
  \bottomrule
\end{longtable}

% =============================================================================
\section{Performance Notes}
% =============================================================================

\begin{itemize}
  \item Trail rendering disabled for $N_P \geq 300$; particle glow also
    disabled at that threshold.
  \item Injection puffs limited to one burst of 5 (one per nozzle) per
    injection event; lifetime $1.42\,\mathrm{s}$.
  \item Encounter cache recomputed at $\approx 5\,\mathrm{Hz}$
    ($O(n^2)$ bound); arc directions recomputed every frame (cheap).
  \item Heatmap finalised once per revolution, not per frame.
  \item KDE orbit band: 200 grid points per population per frame;
    $\approx 200n$ exponential evaluations, negligible at typical
    particle counts.
  \item Hint system: context snapshot built at most once per $12\,\mathrm{s}$
    (global cooldown), not every frame.
  \item Timeline snapshots: one per $10\,\mathrm{s}$ wall-clock time, max 20
    retained. Particle trails and diagnostic flags excluded. Total buffer
    memory $\ll 20\,\mathrm{MB}$, snapshots are pruned when necessary
    to stay within this limit.
\end{itemize}

% =============================================================================
\appendix
% =============================================================================

\section{Tuning Constants Reference}
\label{app:tuning}

\begin{longtable}{@{}lrl@{}}
  \toprule
  Key & Default & Description \\
  \midrule
  \multicolumn{3}{l}{\textit{Drum}} \\
  \texttt{drum.slipRate}         & 0.7      & Omega tracking lag ($\mathrm{s^{-1}}$) \\
  \texttt{drum.omegaDecay}       & 1/30     & Omega target decay ($\mathrm{s^{-1}}$) \\
  \texttt{drum.swipeGain}        & 0.003    & Swipe-to-omega conversion \\
  \texttt{drum.omegaMax}         & $3\pi$   & Max angular velocity ($\mathrm{rad\,s^{-1}}$) \\
  \midrule
  \multicolumn{3}{l}{\textit{Highlight zone}} \\
  \texttt{highlight.cx}          & 50       & Zone centre $x$ (du) \\
  \texttt{highlight.radius}      & 50       & Zone radius (du) \\
  \midrule
  \multicolumn{3}{l}{\textit{Particles}} \\
  \texttt{particle.settleFloor}  & 2.0      & Minimum $\vt$ (cm/s) \\
  \texttt{particle.sizeRefVt}    & 30       & Reference $\vt$ for size scaling \\
  \texttt{particle.collisionR}   & 0.5      & Collision radius (du) \\
  \texttt{particle.showVtProjection} & 0    & 0=off, 1=predicted, 2=live KDE band \\
  \midrule
  \multicolumn{3}{l}{\textit{Aggregates}} \\
  \texttt{aggregate.minLevitated}    & 30   & Min levitated for formation \\
  \texttt{aggregate.mergeCount}      & 10   & Monomers per aggregate \\
  \texttt{aggregate.initialHoldRevs} & 3    & Hold revs before first aggregate \\
  \texttt{aggregate.subseqHoldRevs}  & 1    & Hold revs for subsequent aggregates \\
  \texttt{aggregate.sizeMult}        & 10   & Visual size multiplier \\
  \texttt{aggregate.mergeDur}        & 0.6  & Merge animation duration (s) \\
  \texttt{aggregate.vtFactor}        & 1.2  & Stokes kick factor \\
  \texttt{aggregate.vtGrowFactor}    & 1.05 & Per-growth-merge kick factor \\
  \texttt{aggregate.brownian}        & 2.0  & Brownian kick ($\mathrm{cm\,s^{-1}}$) \\
  \texttt{aggregate.restoreK}        & 0.5  & Restoring spring ($\mathrm{s^{-1}}$) \\
  \texttt{aggregate.flashOrbit}      & false & Enable orbit flash on new aggregate \\
  \texttt{aggregate.growFlashDur}    & 2.5  & Duration of growth orbit flash (s) \\
  \midrule
  \multicolumn{3}{l}{\textit{Pebbles (egg)}} \\
  \texttt{egg.nCrit}             & 7    & Aggregates required for pebble merge \\
  \texttt{egg.holdTarget}        & 3    & Hold revs for first pebble \\
  \texttt{egg.holdSubseq}        & 3    & Hold revs for subsequent pebbles \\
  \texttt{egg.maxBalls}          & 30   & Maximum simultaneous pebbles \\
  \texttt{egg.minSpread}         & 0.30 & Min measured aggregate spread for pebble formation \\
  \midrule
  \multicolumn{3}{l}{\textit{Golden balls}} \\
  \texttt{ball.gravity}          & 200  & Gravitational acceleration ($\mathrm{du\,s^{-2}}$) \\
  \texttt{ball.radius}           & 2.5  & Physical radius (du) \\
  \texttt{ball.wallE}            & 0.85 & Wall restitution coefficient \\
  \texttt{ball.wallFriction}     & 0.92 & Wall tangential retention \\
  \texttt{ball.contactDrag}      & 0.7  & Wall-contact drag ($\mathrm{s^{-1}}$) \\
  \texttt{ball.ballE}            & 0.90 & Ball--ball restitution \\
  \texttt{ball.bumpStrength}     & 0.1  & Bump impulse fraction \\
  \texttt{ball.settleVel}        & 0.6  & Settled speed threshold ($\mathrm{du\,s^{-1}}$) \\
  \midrule
  \multicolumn{3}{l}{\textit{Globes}} \\
  \texttt{globe.nCrit}           & 3      & Pebbles required per globe \\
  \texttt{globe.radius}          & 12.5   & Globe physical radius (du) \\
  \texttt{globe.mergeDur}        & 1.2    & Merge animation duration (s) \\
  \texttt{globe.repulsion}       & 15000  & Globe--globe repulsion constant \\
  \texttt{globe.limit}           & 7      & Maximum simultaneous globes \\
  \midrule
  \multicolumn{3}{l}{\textit{Solar system}} \\
  \texttt{solar.inclCos}         & 0.3    & $\cos i$ of orbital plane tilt \\
  \texttt{solar.omegaBase}       & 3.006  & Keplerian base rate ($\mathrm{rad\,s^{-1}}$) \\
  \texttt{solar.maxR}            & 80     & Max orbit radius (du) \\
  \texttt{solar.persp}           & 0.10   & Perspective size variation \\
  \midrule
  \multicolumn{3}{l}{\textit{Collection tray}} \\
  \texttt{tray.totalAngle}       & $4\pi/3$ & Total insertion sweep ($240^{\circ}$) \\
  \texttt{tray.thickness}        & 3.0    & Blade visual half-thickness (du) \\
  \texttt{tray.decelStart}       & 0.667  & Progress fraction at which braking begins \\
  \texttt{tray.brakeDur}         & 0.8    & Vestigial; not used at runtime \\
  \midrule
  \multicolumn{3}{l}{\textit{Encounter overlay}} \\
  \texttt{encounters.showPairs}  & 10     & Total pairs tracked and drawn \\
  \texttt{encounters.arcPairs}   & 2      & Top-$N$ pairs that show orbit arcs \\
  \midrule
  \multicolumn{3}{l}{\textit{Timeline}} \\
  \texttt{MAX\_SNAPS}            & 20     & Rolling buffer capacity (snapshots) \\
  \texttt{SNAP\_INTERVAL}        & 10     & Wall-clock seconds between snapshots \\
  \midrule
  \multicolumn{3}{l}{\textit{Audio}} \\
  \texttt{audio.masterGain}      & 0.6    & Master volume (0--1) \\
  \texttt{audio.millBlade}       & 0.085  & Mill blade level \\
  \texttt{audio.millHiss}        & 0.10   & Mill hiss level \\
  \texttt{audio.millGrind}       & 0.85   & Mill grind level \\
  \texttt{audio.millThud}        & 0.20   & Mill thud level \\
  \texttt{audio.millBody}        & 0.035  & Mill body level \\
  \texttt{audio.aggMerge}        & 0.4    & Aggregate merge sound level \\
  \bottomrule
\end{longtable}

\end{document}
%%% Local Variables:
%%% mode: LaTeX
%%% TeX-master: t
%%% End:
